\section{The stopped compensated process}\label{sec:process}

Fix \(n\ge2\) and put \(B=n^d\), \(M_n=n^\gamma\), and
\(\epsilon_n=n^{-\zeta}\). A point is active, selected, geometrically
dead, or positively retired. Initially every point is active, its weight
is one, and \(S=\varnothing\). Selection is permanent. Dead and retired
points never return to the active set \(Q_t\).

\subsection{Line variables, rates and flows}
Every designated line has occupancy \(N_\ell=|S\cap\ell|\) and clock
\(k_\ell\). For \(N=0,1\) set
\[
 a_N(k)=\frac{2-N}{2-k},\qquad b_0=\frac12,\quad b_1=1.
\]
Every nondesignated line has a credit \(r_\ell\), initially one.
It changes only while the line contains exactly one selected point.
Before the global stopping time \(\sigma\), define
\begin{align}
 \lambda_p(t)&=\ind_{p\in Q_t}\frac{m_t(p)w_p(t)}n,
 &\Lambda_\ell&=\sum_{p\in\ell}\lambda_p,\label{eq:rates}\\
 \dot k_\ell&=\Lambda_\ell/a_{N_\ell}(k_\ell)
       &&(N_\ell=0,1),\label{eq:clockflow}\\
 \dot r_\ell&=-r_\ell\Lambda_\ell
       &&(N_\ell=1,\ \ell\text{ nondesignated}).\label{eq:tagflow}
\end{align}
The other clock and line-credit derivatives are zero. For active \(p\),
put
\[
 H_p=\sum_{q\in Q_t\setminus\{p\}}\lambda_q
 \ind_{\ell(p,q)\notin\Gamma,\ |S\cap\ell(p,q)|=1}.
\]
Here \(\ell\notin\Gamma\) means that its direction is nondesignated.
The weight flow is
\begin{equation}\label{eq:weightflow}
 \frac{\dot w_p}{w_p}=H_p+\lambda_p+
 \sum_{\substack{\ell\ni p\\\ell\in\Gamma}}
 (\Lambda_\ell-\lambda_p)b_{N_\ell}.
\end{equation}
The term \(H_p\) counts each actual cause \(q\) once, not once for
each possible witness. All rates are set to zero after \(\sigma\).

At a birth of \(q\), use the pre-birth state to choose the branch,
then perform these operations in order:
\begin{enumerate}
\item Select \(q\) and set its weight to zero.
\item Kill every active point that now lies on a line with two selected
points, setting its weight to zero.
\item On each designated line having its first birth, halve the weights
of the remaining active points.
\item Update occupancies and the activation flags of ordinary lines.
The values of \(k_\ell\) and \(r_\ell\) are not reset.
\item Process any newly reached safety boundary and positive retirement.
\end{enumerate}
For a fixed other point \(p\), at most one line through \(p,q\) can
cause a halving. Geometric death takes priority over positive saving.

\subsection{Safety boundaries and the observation convention}
A point retires positively at its first weight-cap hit \(w_p=B\),
at a hit \(r_\ell=\rho\) on an ordinary line through it, or at its
deadline \(\tau_p\). Its current positive weight is saved. Unused
points retire at time zero with saved weight one.

Two additional variables define global guards. For distinct active
\(a,b\), let \(\kappa_{ab}\) be the sum of rates of birth causes
\(q\notin\ell(a,b)\) that would change both weights by a factor less
than one. Set \(\dot\eta_{ab}=-\kappa_{ab}\eta_{ab}\), initially
\(\eta_{ab}=1\), and freeze this variable when either endpoint leaves
the active state. For a designated line \(\ell\) and an active point
\(q\notin\ell\), define the virtual deletion mass
\begin{equation}\label{eq:virtual}
 \mathcal V_{\ell,q}=
 \sum_{y\in\ell\cap Q_t}\frac{w_y}{a_{N_\ell}(k_\ell)}
 \ind_{\ell(q,y)\notin\Gamma,\ |S\cap\ell(q,y)|=1}.
\end{equation}
Globally stop at the first hit of \(k_\ell=K\),
\(\eta_{ab}=1-\epsilon_n\), or \(\mathcal V_{\ell,q}=M_n\), including
an exceedance produced by a jump. Save the positive active credits
and freeze the process. If no guard is hit, \(\sigma=\infty\);
the deterministic terminal time \(T\) is not itself a failure event.

When a mixed observation is saved at a root's positive boundary, take
its value after geometric updates but before removing auxiliary roots
because of simultaneous positive retirement. At a continuous hit take
the boundary values before changing the flags. This convention prevents
a tag witness from disappearing at the very instant it is observed.
It does not alter the selected set or its birth rates. Auxiliary roots
that retired strictly earlier contribute zero.

\begin{proposition}[Construction and legality]\label{prop:construction}
For fixed \(n\), the stopped process is well-defined on \([0,T]\),
is adapted with right-continuous paths and left limits, and always
produces a no-three-in-line set of at most \(2n\) points.
\end{proposition}
\begin{proof}
Give each point a Poisson candidate clock of rate \(B/n\), together
with independent uniform marks. At a candidate arrival, accept the
birth if its mark is at most \(m_t(p)w_p(t-)/B\) and the point is
eligible in the pre-arrival state. The total candidate rate is \(nB\),
so there are finitely many arrivals on \([0,T]\) almost surely.
Use the filtration revealing only arrivals and marks up to the current
time, with the usual completion and right-continuous modification.

Between arrivals and mode changes, solve the finite-dimensional ODE.
Within the stopped region the denominators are separated from zero:
\(k\le K<2\), \(r\ge\rho>0\), and \(\eta\ge1-\epsilon_n>0\).
Using \(w\), rather than \(\log w\), as a coordinate, each mode has a
locally Lipschitz vector field. Extend it outside a compact
neighborhood with a cutoff if necessary. The resulting flow is unique
and depends continuously on its initial state. First hits of its
closed guards are measurable: on any bounded time interval, the
minimum distance of the continuous flow to a closed guard detects a
hit. Composition at the preceding stopping time preserves adaptation.

Every positive retirement occurs at most once per point. Occupancies
and line activation flags change only finitely often, deadlines and
phases are finite, and a global stop occurs once. Resolve all
simultaneous guards together and mark them as processed. Together with
the finite candidate arrivals, this rules out an accumulation of
updates. The construction gives the claimed path regularity.

Immediately before every birth the chosen point is active, hence is
not on a line with two selected points. Adding it preserves legality;
the subsequent deaths maintain this invariant. The row bound gives
\(|S|\le2n\).
\end{proof}

For compensators, evaluate the intensity in the pre-jump state and use
\(\ind_{t\le\sigma}\). The birth that triggers a global stop is counted.
In time integrals one may instead use \(\ind_{t<\sigma}\), since the
two differ on at most one time per path. Thinning gives, for a
nonnegative predictable observation \(g\),
\begin{equation}\label{eq:compensator}
 \E\sum_{\text{births }(p,t)}g_p(t)
 =\E\int_0^T\sum_p\lambda_p(t)g_p(t)\dd t.
\end{equation}
For simple predictable functions this is independent increments and
the uniform-mark rule; monotone approximation gives the general
identity. Bounded signed observations follow by subtraction. All
fixed-\(n\) rates and stopped observations used below are integrable.

\section{Saved credit and the exact output identity}\label{sec:loss}

Define \(Z_p\) to equal \(w_p\) while \(p\) is active, zero after
selection or geometric death, and its saved value after positive
retirement or a global stop. Let \(C_p\) be the birth-eligible credit,
including the global-stop indicator, and let \(V_p\) be saved credit.
Then \(Z_p=C_p+V_p\) and \(0\le Z_p\le B\).

\begin{lemma}[One-point compensation]\label{lem:credit}
For each point \(p\), \(Z_p\) is a bounded martingale and
\(\E Z_p(t)=1\).
\end{lemma}
\begin{proof}
While \(p\) is active, let \(J_p(q)\in\{0,1/2,1\}\) be its
multiplier at a birth of \(q\). Selection or deletion gives zero,
a first designated-line birth gives \(1/2\), and otherwise it is one.
The flow \eqref{eq:weightflow} is precisely
\[
 \dot w_p=w_p\sum_q\lambda_q(1-J_p(q)).
\]
By \eqref{eq:compensator}, the jump drift is its negative. Positive
saving changes no value; subsequent drift is zero. A coincident
geometric death has already been assigned multiplier zero. Thus the
compensated drift vanishes, and boundedness makes the result a true
martingale. Its initial value, including immediate saving of unused
points, is one.
\end{proof}

\begin{lemma}[Point law and birth intensities]\label{lem:pointlaw}
The formula \(\dd\Prob^{(p)}=Z_p(T)\dd\Prob\) defines a probability
measure. At any bounded stopping time \(\theta\le T\),
\(\E^{(p)}X=\E[Z_p(\theta)X]\) for nonnegative
\(\F_\theta\)-measurable \(X\). Before the root's saving boundary,
the birth intensity of \(q\) under this law is \(J_p(q)\lambda_q\).
\end{lemma}
\begin{proof}
The martingale identity gives
\(\E[Z_p(T)\mid\F_\theta]=Z_p(\theta)\).
One can obtain the stopping-time version from finite-valued dyadic
approximations from above, right continuity and bounded convergence.
For the intensity statement, enumerate the births, condition
\(Z_p(T)\) on each birth's post-jump sigma field, substitute
\(Z_p(t)=J_p(q)Z_p(t-)\), and then apply
\eqref{eq:compensator}. Tonelli justifies the nonnegative birth sum.
On time integrals \(Z_p(t-)=Z_p(t)\) almost everywhere.
\end{proof}

The changed law forbids a birth that kills the root, but does not
delete the self-rate term \(\lambda_p\) from the original deterministic
weight ODE. Losing this distinction changes the compensation identity.

\begin{proposition}[Exact remaining-exposure loss]\label{prop:loss}
With \(E_p\) and \(\theta_p\) as in Section~\ref{sec:overview},
\eqref{eq:previewloss} holds for every fixed \(n\). Therefore
\[
 \sup_p\Prob^{(p)}(E_p)\le\varepsilon
 \quad\Longrightarrow\quad
 \E|S_n|\ge(1-\varepsilon)n A_n.
\]
\end{proposition}
\begin{proof}
Apply \eqref{eq:compensator} to the birth count and use \(\E Z_p(t)=1\):
\[
 \E|S_n|=\frac1n\sum_p\int_0^T m_t(p)\E C_p(t)\dd t
 =nA_n-\frac1n\sum_p\int_0^T m_t(p)\E V_p(t)\dd t.
\]
On a path with positive saving,
\(V_p(t)=Z_p(T)\ind_{\theta_p\le t}\). Tonelli changes its integral
to \(Z_p(T)\int_{\theta_p}^T m_t(p)\dd t\).
Deadline saving has no remaining exposure. A path dying without
positive saving has zero terminal credit and may be assigned
\(\theta_p=T\). The density identity proves \eqref{eq:previewloss}.
Finally the remaining exposure is between zero and the deterministic
total \(F_{\typ(p),T}(p/n)\); sum the probability bound. Events need
not be independent or disjoint.
\end{proof}

\begin{lemma}[Product formula and a positive lower bound]\label{lem:product}
While \(p\) is active,
\begin{equation}\label{eq:product}
 w_p(t)=
 \exp\left\{\int_0^t H_p(u)\dd u+
 \int_0^t\lambda_p(u)\left(1-
       \sum_{\substack{\ell\ni p\\\ell\in\Gamma}}b_{N_\ell}\right)\dd u\right\}
 \prod_{\substack{\ell\ni p\\\ell\in\Gamma}}a_{N_\ell}(k_\ell).
\end{equation}
For large enough \(n\), \(w_p\ge2^{-g-1}\) uniformly while active.
The absolute value of the self-correction integral is at most
\(\max(1,g-1)TB/n=o(1)\).
\end{lemma}
\begin{proof}
Between births, \(\mathrm d\log a_N(k)/\mathrm dt=\Lambda_\ell b_N\).
At a first designated-line birth \(a_1/a_0=1/2\), exactly matching
the surviving point's halving. Subtract these factors from
\eqref{eq:weightflow}; all designated-line terms cancel except the
displayed self-correction. Initially every factor is one. Since
\(H_p\ge0\), \(a_0\ge1\), \(a_1\ge1/2\), and \(\lambda_p\le B/n\),
the asserted bounds follow.
\end{proof}

